\documentclass[11pt]{amsart}
\usepackage[margin=1.05in]{geometry}
\usepackage{amsmath,amssymb}
\usepackage{graphicx}
\usepackage[colorlinks=true,linkcolor=blue,citecolor=blue,urlcolor=blue]{hyperref}

\newcommand{\R}{\mathbb R}
\newcommand{\T}{\mathsf T}
\newcommand{\norm}[1]{\left\lVert #1\right\rVert}

\title{A Counterexample to a Real-\(p\) Mixed Matrix Norm Extension}
\author{fa\_banach\_001 counterexample packet}
\date{2026-07-03}

\begin{document}
\maketitle

\section*{Status}

This packet records a counterexample to a question in Sah--Sawhney--Stoner--Zhao,
\emph{A reverse Sidorenko inequality}, arXiv:1809.09462.  In Section 3.1,
after Lemma 3.1, they define
\[
  \norm{M}_{L_{p,q}}
  =
  \left(\sum_i \left(\sum_j |m_{ij}|^p\right)^{q/p}\right)^{1/q}
\]
and ask whether, for nonnegative matrices \(A,B\), the inequality
\[
  \norm{A^\T B}_{L_{p,q}}^2
  \le
  \norm{A^\T A}_{L_{q,q}}\,\norm{B^\T B}_{L_{p,p}}
\tag{1}
\]
extends from positive integer \(p\) to all real \(1\le p\le q\).

\begin{center}
\includegraphics[width=0.94\linewidth]{figures/open_problem_crop.png}
\end{center}

\section*{Idea of the Proof}

When \(p=q\), the mixed norm is just the entrywise \(\ell_p\)-norm.  For
integer \(p\), the kernel \(\langle x,y\rangle^p\) is positive semidefinite
because it is the inner product of tensor powers.  The natural place to look
for a non-integer obstruction is therefore a small positive-measure
configuration for the kernel \(\langle x,y\rangle^p\) on the positive cone.
Two directions in \(\R_+^2\), with reciprocal weights, already break the
Cauchy--Schwarz-type inequality at \(p=3/2\).

\section*{Counterexample}

Set
\[
  p=q=\frac32,\qquad t=\frac14,\qquad r=1-t^p=\frac78,\qquad
  \alpha=r^{1/p}=r^{2/3}.
\]
Let
\[
  A=\begin{pmatrix}
      0 & 1\\
      \alpha & t
    \end{pmatrix},
  \qquad
  B=\begin{pmatrix}
      t & \alpha\\
      1 & 0
    \end{pmatrix}.
\]
These are nonnegative real \(2\times2\) matrices.

\section*{Proof}

Since \(p=q\), write \(\norm{M}_p=(\sum_{ij}|m_{ij}|^p)^{1/p}\).  The proposed
inequality (1) becomes
\[
  \norm{A^\T B}_p^2 \le \norm{A^\T A}_p\,\norm{B^\T B}_p.
\]
The columns of \(A\) are \(\alpha e_2\) and \((1,t)\), while the columns of
\(B\) are \((t,1)\) and \(\alpha e_1\).  Thus
\[
  \norm{A^\T B}_p^p
  =
  (\alpha)^p + (2t)^p + (\alpha)^p
  =
  2r+(2t)^p=:S.
\]
Similarly,
\[
  \norm{A^\T A}_p^p
  =
  \alpha^{2p}+2(\alpha t)^p+(1+t^2)^p
  =
  r^2+2rt^p+(1+t^2)^p=:X.
\]
By symmetry \(\norm{B^\T B}_p^p=X\).  Therefore (1) is equivalent to
\[
  S^{2/p}\le X^{2/p},
\]
or simply \(S\le X\).

For \(t=1/4\) and \(p=3/2\),
\[
  S=2\cdot\frac78+\left(\frac12\right)^{3/2}
   =\frac{7+\sqrt2}{4},
\]
whereas
\[
  X=\left(\frac78\right)^2
     +2\cdot\frac78\cdot\frac18
     +\left(\frac{17}{16}\right)^{3/2}
   =\frac{63+17\sqrt{17}}{64}.
\]
It remains only to compare these two numbers.  Multiplying by \(64\), the
claim \(S>X\) is
\[
  112+16\sqrt2 > 63+17\sqrt{17},
\]
or
\[
  49+16\sqrt2>17\sqrt{17}.
\]
This holds because
\[
  (49+16\sqrt2)^2-(17\sqrt{17})^2
  =
  1568\sqrt2-2000>0,
\]
and the final inequality follows from \(\sqrt2>125/98\).  Hence \(S>X\), so
\[
  \norm{A^\T B}_{L_{3/2,3/2}}^2
  >
  \norm{A^\T A}_{L_{3/2,3/2}}\,
  \norm{B^\T B}_{L_{3/2,3/2}}.
\]
This disproves the proposed real-\(p\) extension.

\section*{Verification Notes}

The script \texttt{code/verify\_counterexample.py} evaluates the same example
numerically and prints the ratio
\[
  \frac{\norm{A^\T B}_{L_{3/2,3/2}}^2}
       {\norm{A^\T A}_{L_{3/2,3/2}}\norm{B^\T B}_{L_{3/2,3/2}}},
\]
which is approximately \(1.0154>1\).  The proof above is exact; the script is
only a check against transcription errors.

\section*{Novelty and Limitations}

Cheap run indexes were searched for arXiv:1809.09462, ``reverse Sidorenko'',
``mixed norm'', and the exact matrix inequality; no prior run packet or
attempt was found.  Local parsed-source search found the exact inequality only
in the source paper and generated target-pool extracts.  Web searches on
2026-07-03 for exact phrases from the inequality and the source sentence found
no visible later answer.  This is a bounded novelty check, not a publication
claim.

The counterexample only disproves the proposed extension to all real \(p\). It
does not contradict the source paper's \(p=1\) lemma, the positive-integer
cases obtained by tensor powers, or any of the reverse Sidorenko theorems.

\section*{Reference}

A. Sah, M. Sawhney, D. Stoner, and Y. Zhao,
\emph{A reverse Sidorenko inequality}, arXiv:1809.09462.

\end{document}
